\section{Open Problems and Future Directions}
\label{sec:open_problems}

The GEM framework presented in this work establishes the topological origin of fundamental couplings through the discrete icosahedral structure of the quantum vacuum. However, several important questions remain open for future investigation.


\subsection{The Hierarchy Problem in Muon g-2 and Future Directions}
\label{subsec:hierarchy_g2}

A significant prediction of the GEM framework is the existence of a topological correction to the muon anomalous magnetic moment, arising from the interaction of the muon with the discrete icosahedral vacuum structure. However, this prediction reveals a hierarchy problem that requires careful analysis.

If the icosahedral symmetry breaking occurs at the Planck scale ($M_{I_h} \sim M_{\text{Planck}} \approx 10^{19}$ GeV), as postulated in Section 2, a naive dimensional estimate yields a correction approximately 34 orders of magnitude smaller than the experimentally observed discrepancy ($\Delta a_\mu \approx 2.5 \times 10^{-9}$).

This significant mismatch is conceptually analogous to the well-known hierarchy problem of the Higgs boson mass in the Standard Model, where the electroweak scale ($M_{EW} \sim 100$ GeV) is unnaturally light compared to the Planck scale without a protective symmetry. Following this analogy, a satisfactory resolution of this GEM hierarchy problem should explain why the effective coupling scale $M_{I_h}^{\text{eff}}$ differs from $M_{\text{Planck}}$ by $\sim 17$ orders of magnitude without introducing an equivalent fine-tuning. 

We identify several potential avenues for future research to address this:
\begin{enumerate}
    \item \textbf{Renormalization Group (RG) Running:} The torsion-fermion coupling may exhibit anomalous dimension running between the Planck scale and the electroweak scale, naturally enhancing the effective coupling at low energies.
    \item \textbf{Light Torsion Modes:} The icosahedral vacuum structure may possess collective excitations or extended topological defects at intermediate energy scales, analogous to the relationship between $\Lambda_{\text{QCD}}$ and hadron masses.
    \item \textbf{Higgs-Mediated Coupling:} The interaction may be suppressed by the Higgs mechanism, requiring a detailed analysis of the torsion condensate effective potential and its coupling to the Higgs sector.
\end{enumerate}

The determination of the precise mechanism and the effective coupling scale $M_{I_h}^{\text{eff}}$ constitutes a primary objective for future work. We note that the current framework successfully provides the \textit{topological origin} of the correction, while its precise quantitative magnitude at low energies remains an open problem deeply connected to the fundamental hierarchies of particle physics.

\subsection{Mass Hierarchy and Renormalization Effects}
\label{subsec:mass_hierarchy}

The qualitative prediction of three fermion generations with increasing masses (Proposition 5.2) captures the correct ordering ($m_e < m_\mu < m_\tau$) but underestimates the mass ratios by several orders of magnitude. This discrepancy indicates that the tree-level approximation $m_k \propto n_k^2$ must be supplemented by radiative corrections and renormalization group effects, which are expected to be substantial given the strong topological coupling. A complete calculation of these effects is beyond the scope of the present work.

\subsection{Experimental Validation at Low Energies}
\label{subsec:experimental_validation}

The macroscopic predictions of the GEM framework (anomalous cooling and inertial mass reduction in water at 16.2 Hz) represent an exploratory hypothesis that requires careful experimental validation. The mechanism by which topological torsion at the Planck scale could manifest as measurable macroscopic effects at room temperature remains to be fully elucidated. Experimental efforts along these lines will provide crucial tests of the framework's predictions in the low-energy regime.